%ceci est un commentaire efficace
%ceci est un deuxieme commentaire efficace
\documentclass{beamer}

    \usepackage[utf8]{inputenc}
    \usepackage[T1]{fontenc}
    \usepackage[french]{babel}
    \usepackage{url}

    \usetheme{Warsaw}

    \logo{%
            \includegraphics[scale=0.06]{logo}
    }

    \title[Coworking-oIP]{Coworking Over IP}
    \author{Outils et astuces \textit{over IP} pour le travail collaboratif}
    \institute{Mathieu (matael) Gaborit | HAUM | Ruche Numérique}
    \date{2013}

\begin{document}

\begin{frame}
\titlepage
\end{frame}

\begin{frame}
\frametitle{Au Menu}
\tableofcontents
\end{frame}

\section{Avant de commencer} % {{{1

\begin{frame}
    \frametitle{HAUM}

    \begin{itemize}
        \item HAckerspace de l'Université du Maine
        \item association de passionné(e)s de \textit{hacking} et d'expérimentations
        \item par extension, association de technophiles :)
        \item Site Web : \url{http://haum.org}
        \item Twitter : \texttt{@haum72}
        \item Mailing-list : \url{http://lists.matael.org/mailman/listinfo/haum_hackerspace}
        \item IRC : \texttt{\#haum @ irc.freenode.org}
    \end{itemize}
\end{frame}

\begin{frame}
    \frametitle{La Ruche Numérique}

    \begin{itemize}
        \item TODO
    \end{itemize}
\end{frame}

\begin{frame}
    \frametitle{Propriétaire ou libre ?} 

    \centering{\Huge{\texttt{<troll></troll>}}}
\end{frame}

\section{Constat}

\begin{frame}
    \frametitle{Constat}

    \begin{itemize}[<+->]
        \item augmentation du nombre de télétravailleurs
        \item important besoin de communication avec le reste du monde
        \item insuffisance des outils \textit{standards}
        \item besoin de travailler en même temps sur un même produit...
    \end{itemize}
\end{frame}

\section{Outils standards} % {{{1

\subsection{Old School} % {{{2

\begin{frame}
    \frametitle{Depuis des outils très \textit{old school}...}

    \pause{}

    \begin{itemize}
        \item téléphone
        \item lettre (oui oui...)
        \item rencontre \textit{IRL} (\textit{in real life})
    \end{itemize}
\end{frame}

\begin{frame}
    \begin{block}{Avantages}
    \begin{itemize}
        \item bien connus et maitrisés
        \item \textit{bullet proof}
        \item canaux priviligiés desservant (presque) tout le monde
    \end{itemize}
    \end{block}

    \pause{}

    \begin{block}{Inconvénients}
    \begin{itemize}
        \item pas toujours adaptés à la réalité actuelle
            \note[item]{se déplacer ou envoyer une lettre pour envoyer un document plutot qu'un mail => :S}
        \item ne permettent pas (ou peu) le travail collaboratif
        \item peuvent engendrer des frais (lettres recommandées, etc...)
            \note[item]{à ceux qui trouvent le mail pas fiable, nous allons y revenir}
    \end{itemize}
    \end{block}

\end{frame}

\subsection{Online} % {{{2

\begin{frame}
    \frametitle{... à du standard \textit{online}....} 

    \begin{itemize}
        \item email
        \item Skype
        \item Twitter/Facebook
    \end{itemize}

    \pause{}

    \begin{block}{Bilan des forces}
        \begin{itemize}
            \item bien connus aussi mais souvent mal maîtrisés
            \item permettent l'envoi à plusieurs
            \item robustes et plus ou moins éprouvés
        \end{itemize}
    \end{block}
\end{frame}

\subsection{Le reste...} % {{{2

\begin{frame}
    \frametitle{... voire de l'inconnu}

    \begin{itemize}
        \item AdobeConnect (salle de réunion over IP), propriétaire
        \item Mumble : alternative libre de chat vocal par salon
            \note[item]{Utilisé par le PP et d'autres groupes notament}
    \end{itemize}

    \pause{}

    \begin{block}{Avis}
        \begin{itemize}
            \item méconnus mais puissants
            \item flexibles et peu couteux
            \item ...
        \end{itemize}

    \end{block}
\end{frame}

\subsection{Des mails...} % {{{2

\begin{frame}
    \frametitle{Mail et sécurité}

    \begin{itemize}
        \item Créé par la RFC822 en 1982
        \item joue le rôle électronique de La Poste physique ... avec tous ses travers.
    \end{itemize}

    \pause{}
    
    La poste est un \textit{intermédiaire de confiance}, le mail aussi, mais on ne le connait pas ni ne pouvons le
    contrôler.

    \pause{}

    \medskip

    Il est important de :

    \begin{itemize}
        \item Signer ses messages (via PGP/GPG par exemple)
        \item Chiffrer les données sensibles
        \item Rester loin des CGU obscures et dangereuses de certains fournisseurs (Google,...)
    \end{itemize}
\end{frame}

\section{Des outils moins courants} % {{{1

\begin{frame}
    \frametitle{Pourquoi toutes ces solutions ?}

    \begin{itemize}
        \item Répondre à des besoins précis
        \item Pour adapter au mieux les outils à son travail et non l'inverse
        \item Pour rester réactif et efficace
    \end{itemize}

    \pause{}

    Quatre exemples pour quatre besoins.
\end{frame}

\subsection{Etherpad(Lite)} % {{{2

\begin{frame}
    \frametitle{Exemple 1 : Besoins}

    \begin{itemize}
        \item rédiger en commun des documents
        \item \textit{brainstorming}
    \end{itemize}

\end{frame}


\begin{frame}
    \frametitle{Exemple 1 : Etherpad(Lite)}

    \begin{itemize}
        \item dispo dans de nombreux points du web (Framasoft, divers Partis Pirates, autres, ...)
        \item permet l'export, la création "d'équipes", etc...
        \item \textbf{Attention} charge maxi variable selon fournisseur
        \item simple à installer sur son propre serveur
    \end{itemize}

    \begin{block}{Implémentations existantes}
        \begin{itemize}
            \item Etherpad: Java
            \item EtherpadLite: Node.js
        \end{itemize}
    \end{block}


\end{frame}

\subsection{Write\LaTeX} % {{{2

\begin{frame}
    \frametitle{Exemple 2 : Besoins}

    \begin{itemize}
        \item rédiger des documents avec une bonne mise en page ou des slides
        \item utiliser \LaTeX{} en collaboration
    \end{itemize}

\end{frame}


\begin{frame}
    \frametitle{Exemple 2 : Write\LaTeX{}}

    \begin{itemize}
        \item une sorte d'Etherpad adapté à LaTeX avec auto-régénération du résultat et modifications partagées
        \item testé et plus qu'approuvé
    \end{itemize}

\end{frame}


\subsection{Git} % {{{2

\begin{frame}
    \frametitle{Exemple 3 : Besoins}

    \begin{itemize}
        \item rédaction de gros documents en collaboratif
        \item code communautaire
        \item gestion de projet
        \item gestion de liste de tâches et attribution
    \end{itemize}

\end{frame}

\begin{frame}
    \frametitle{Exemple 3 : Serveur Git !}

    \begin{itemize}
        \item Github : payant pour un compte privé
        \item BitBucket : compte privés/publics gratuit
        \item Gitolite sur un serveur perso
    \end{itemize}

    \pause

    \begin{block}{Possibilités}

        \begin{itemize}
            \item un suivi des modifications
            \item gestion des tickets
            \item attribution des tâches
        \end{itemize}
    \end{block}

\end{frame}

\begin{frame}
    \frametitle{Gestion de projet pure}

    \begin{itemize}
        \item Redmine
        \item Indefero
        \item Trac
        \item PivotalTracker (privé, payant)
    \end{itemize}

    \pause

    \begin{block}{Autres infos}
        \begin{itemize}
            \item Github a une appli Androïd et des clients Windows et Mac
            \item Tortoise-git est une appli git pour Windows
        \end{itemize}
    \end{block}
\end{frame}

\subsection{OwnCloud} % {{{2

\begin{frame}
    \frametitle{Exemple 4 : Besoins}

    \begin{itemize}
        \item partage de fichiers sur plusieurs machines avec plusieurs personnes
        \item agenda commun
        \item serveur perso
    \end{itemize}
\end{frame}

\begin{frame}
    \frametitle{Exemple 4 : OwnCloud}

    \begin{itemize}
        \item enfin un cloud libre et perso
        \item supporté par une vraie communauté
        \item utilisation sur serveur communautaire
    \end{itemize}
\end{frame}

\section{Gare aux arnaques} % {{{1

\begin{frame}
    \frametitle{Gare Aux Arnaques}

    Avant de choisir définitivement un service, pensez à lire les CGU.

    \pause{}

    Des sites peuvent aider :

    \centering{\url{http://tosdr.org/}}

\end{frame}

\begin{frame}
    \frametitle{Gare Aux Arnaques}

    \begin{block}{Dropbox}

        Dropbox et compagnie gardent souvent des droits sur les contenus déposés.

        Elles restent malheureusement très utilisées par certaines entreprises....

    \end{block}

\end{frame}

\begin{frame}
    \frametitle{Gare Aux Arnaques}

    \begin{block}{Google}

        Connu pour ses pratiques douteuses (suppression de compte sans préavis).

    \end{block}
    
    \pause{}

    \begin{block}{Apple iCloud}

        C'est globalement une mauvaise idée : Apple Inc se réserve un contrôle total des données déposées

    \end{block}


\end{frame}
\begin{frame}
    \frametitle{Gare Aux Arnaques}


    \begin{block}{Mega}

        Tout nouveau sur le marché, Mega est intéressant, attention à l'article 8 des CGU...

        \centering{\small{\textit{8. Our service may automatically delete a piece of data you upload or give someone else access to where
            it determines that that data is an exact duplicate of original data already on our service. In that case,
you will access that original data.}}}

    \end{block}
\end{frame}
% }}}

\begin{frame}
    
    \frametitle{Remerciements}

    \begin{itemize}
        \item La Ruche Numérique et la CCI
        \item LinuXMaine
        \item Le HAUM
        \item Les relecteurs de ces slides et du plan
    \end{itemize}

    \pause{}

    \centering{\Huge{Vous ... pour votre écoute}}

\end{frame}
\end{document}

